Com base na Resolução CNE/CES nº 5/2016, fundamentada no Parecer CNE/ CES nº 136/2012, que estabelece os conteúdos curriculares para os cursos da área de Computação (DCN Computação), apresentam-se a seguir os componentes curriculares do curso que contemplam esses conteúdos de forma transversal e integrada.

No referido parecer, estão definidos os \textbf{conteúdos tecnológicos e básicos comuns aos cursos de Bacharelado e Licenciatura} da área, abrangendo fundamentos teóricos da computação, desenvolvimento de software, infraestrutura e sistemas computacionais, além de aspectos éticos, sociais e econômicos relacionados à profissão.

O documento também especifica os \textbf{conteúdos próprios da Engenharia de Computação}, voltados ao domínio de hardware, sistemas embarcados, eletrônica, controle e automação, consolidando a interface entre a Computação e as Engenharias.

A seleção dos conteúdos apresentados nos quadros seguintes reflete o entendimento do NDE e do Colegiado do Curso quanto à aderência entre os componentes curriculares e os eixos formativos estabelecidos pelas Diretrizes Nacionais, garantindo a formação abrangente e atualizada prevista na norma.

\begin{longtblr}[
        theme = ecp,
        caption = {Distribuição dos Conteúdos Curriculares Referentes à Formação Tecnológica e Básica para todos os Cursos de Bacharelado e de Licenciatura (DCN Computação)},
        label = {tab:disciplinas_ref_dcn_base},
        remark{\small \textbf{Referência}} = {\small Parecer CNE/CES Nº 136/2012, item 3.1.}]{
    colspec = {Q[c,m,wd=20mm]Q[l,m,wd=50mm]Q[l,m,wd=25mm]Q[l,m,wd=40mm]},
    rowhead = 1,
    row{odd} = {bg=CinzaClaro},
    row{1} = {bg=AzulEscuro, fg=white},
    row{60} = {bg=AzulClaro, fg=white},
    cells  = {font=\fontsize{10pt}{12pt}\selectfont},
    }
        \textbf{Período} & \textbf{Componente} & \textbf{Código} & \textbf{Conteúdos}\\
        1\textordmasculine & Cálculo Diferencial e Integral I & FAMAT31011 & {\scriptsize matemática do contínuo} \\
    1\textordmasculine & Geometria Analítica & FAMAT31021 & {\scriptsize matemática do contínuo} \\
    1\textordmasculine & Introdução à Engenharia de Computação & FEELT31103 & {\scriptsize computação e sociedade; empreendedorismo; filosofia; meio ambiente; metodologia cientifica; ética e legislação} \\
    1\textordmasculine & Introdução à Prototipagem Eletrônica & FEELT!IPE & {\scriptsize automação; sistemas embarcados} \\
    1\textordmasculine & Lógica e Matemática Discreta para Computação & FEELT!LMD & {\scriptsize análise combinatória; lógica; matemática discreta; métodos formais} \\
    1\textordmasculine & Programação Script & FEELT!PS & {\scriptsize programação} \\
    2\textordmasculine & Álgebra Linear & FAMAT31022 & {\scriptsize estruturas algébricas; matemática do contínuo} \\
    2\textordmasculine & Cálculo Diferencial e Integral II & FAMAT31012 & {\scriptsize matemática do contínuo} \\
    2\textordmasculine & Física Básica: Mecânica & INFIS39031 & {\scriptsize matemática do contínuo} \\
    2\textordmasculine & Experimental de Física Básica: Mecânica & INFIS39032 & {\scriptsize matemática do contínuo} \\
    2\textordmasculine & Metrologia & FEELT31204 & {\scriptsize automação; matemática do contínuo; meio ambiente} \\
    2\textordmasculine & Programação Procedimental & FEELT!PP & {\scriptsize programação} \\
    2\textordmasculine & Sistemas Digitais & FEELT31409 & {\scriptsize circuitos digitais} \\
    2\textordmasculine & Experimental de Sistemas Digitais & FEELT31410 & {\scriptsize circuitos digitais} \\
    3\textordmasculine & Cálculo Diferencial e Integral III & FAMAT31013 & {\scriptsize matemática do contínuo} \\
    3\textordmasculine & Elementos de Sistemas Computacionais & FEELT!ESC & {\scriptsize arquitetura e organização de computadores; circuitos digitais; compiladores} \\
    3\textordmasculine & Estrutura de Dados & FEELT!ED & {\scriptsize abstração e estruturas de dados; programação} \\
    3\textordmasculine & Física Básica: Eletricidade e Magnetismo & INFIS39033 & {\scriptsize matemática do contínuo} \\
    3\textordmasculine & Experimental de Física Básica: Eletricidade e Magnetismo & INFIS39034 & {\scriptsize matemática do contínuo} \\
    3\textordmasculine & Inteligência Artificial, Ética e Sociedade & FEELT!IAES & {\scriptsize computação e sociedade; filosofia; inteligência artificial e computacional; meio ambiente; ética e legislação} \\
    3\textordmasculine & Programação Funcional & FEELT!PF & {\scriptsize estruturas algébricas; fundamentos de linguagens; novos paradigmas de computação; programação} \\
    3\textordmasculine & Programação Orientada a Objetos & FEELT!POO & {\scriptsize engenharia de software; programação} \\
    4\textordmasculine & Análise de Algoritmos & FEELT!AA & {\scriptsize abstração e estruturas de dados; algoritmos e complexidade; teoria dos grafos} \\
    4\textordmasculine & Arquitetura e Organização de Computadores & FEELT!AOC & {\scriptsize arquitetura e organização de computadores} \\
    4\textordmasculine & Atividades Curriculares de Extensão I & ACE1 & {\scriptsize computação e sociedade} \\
    4\textordmasculine & Banco de Dados & FEELT!BD & {\scriptsize banco de dados; programação} \\
    4\textordmasculine & Estatística & FAMAT31033 & {\scriptsize probabilidade e estatística} \\
    4\textordmasculine & Métodos Matemáticos & FAMAT31031 & {\scriptsize matemática do contínuo} \\
    5\textordmasculine & Atividades Curriculares de Extensão II & ACE2 & {\scriptsize computação e sociedade} \\
    5\textordmasculine & Cálculo Numérico & FAMAT31032 & {\scriptsize matemática do contínuo; modelagem computacional} \\
    5\textordmasculine & Empreendedorismo e Inovação & FAGEN39401 & {\scriptsize empreendedorismo; fundamentos de administração} \\
    5\textordmasculine & Engenharia de Software & FEELT!ESOF & {\scriptsize análise, especificação, verificação e testes de sistemas; engenharia de software} \\
    5\textordmasculine & Fundamentos da Inteligência Artificial & FEELT!FIA & {\scriptsize inteligência artificial e computacional; lógica; programação} \\
    5\textordmasculine & Sistemas Operacionais & FEELT!SO & {\scriptsize sistemas operacionais} \\
    6\textordmasculine & Aprendizagem de Máquina & FEELT!AMAQ & {\scriptsize inteligência artificial e computacional; modelagem computacional; probabilidade e estatística} \\
    6\textordmasculine & Atividades Curriculares de Extensão III & ACE3 & {\scriptsize computação e sociedade} \\
    6\textordmasculine & Processamento Digital de Sinais & FEELT31628 & {\scriptsize matemática do contínuo; modelagem computacional} \\
    6\textordmasculine & Redes de Comunicações I & FEELT31724 & {\scriptsize avaliação de desempenho; redes de computadores} \\
    6\textordmasculine & Sistemas Embarcados I & FEELT31523 & {\scriptsize automação; sistemas de tempo real; sistemas embarcados; sistemas operacionais} \\
    6\textordmasculine & Teoria da Computação & FEELT!TCOMP & {\scriptsize algoritmos e complexidade; linguagens formais e autômatos; métodos formais; teoria da computação} \\
    7\textordmasculine & Aprendizagem Profunda e Modelos Generativos & FEELT!APMG & {\scriptsize inteligência artificial e computacional; modelagem computacional; processamento de imagens} \\
    7\textordmasculine & Atividades Curriculares de Extensão IV & ACE4 & {\scriptsize computação e sociedade} \\
    7\textordmasculine & Redes de Comunicações II & FEELT31831 & {\scriptsize avaliação de desempenho; redes de computadores} \\
    7\textordmasculine & Segurança de Sistemas Computacionais & FEELT!SEG & {\scriptsize redes de computadores; segurança; sistemas operacionais} \\
    7\textordmasculine & Sistemas Computacionais em Tempo Real & FEELT!SCTR & {\scriptsize avaliação de desempenho; dependabilidade; sistemas de tempo real; sistemas operacionais} \\
    7\textordmasculine & Sistemas e Controle & FEELT!SCON & {\scriptsize automação; modelagem computacional} \\
    8\textordmasculine & Arquitetura de Software Aplicada & FEELT!ASA & {\scriptsize engenharia de software; processamento distribuído; processamento paralelo} \\
    8\textordmasculine & Engenharia Econômica & IERI39601 & {\scriptsize fundamentos de administração; fundamentos de economia} \\
    8\textordmasculine & Engenharia de Contexto em IA Generativa & FEELT!ECIA & {\scriptsize computação e sociedade; inteligência artificial e computacional; interação humano-computador; ética e legislação} \\
    8\textordmasculine & IA Aplicada à Cibersegurança & FEELT!IACS & {\scriptsize inteligência artificial e computacional; redes de computadores; segurança} \\
    8\textordmasculine & Memorial de Atividades Curriculares de Extensão & FEELT!MACE & {\scriptsize computação e sociedade; metodologia cientifica} \\
    8\textordmasculine & Sistemas Distribuídos & FEELT!SD & {\scriptsize avaliação de desempenho; banco de dados; dependabilidade; processamento distribuído} \\
    Optativa & Compiladores & FACOM39043 & {\scriptsize compiladores} \\
    Optativa & Computação Gráfica RV-RA & FEELT31721 & {\scriptsize computação gráfica; realidade virtual} \\
    Optativa & Multimídia & GBC213 & {\scriptsize multimídia} \\
    Optativa & Otimização & FACOM31601 & {\scriptsize pesquisa operacional e otimização} \\
    Optativa & Pesquisa Operacional & FAGEN31703 & {\scriptsize pesquisa operacional e otimização} \\
    Optativa & Robótica & FEELT31722 & {\scriptsize robótica} \\*

    \end{longtblr}
    

\begin{longtblr}[
        theme = ecp,
        caption = {Distribuição dos Conteúdos Curriculares Referentes à Formação Tecnológica e Básica dos Cursos de Bacharelado em Engenharia de Computação (DCN Computação)},
        label = {tab:disciplinas_ref_dcn_ecp},
        remark{\small \textbf{Referência}} = {\small Parecer CNE/CES Nº 136/2012, item 3.3.}]{
    colspec = {Q[c,m,wd=20mm]Q[l,m,wd=50mm]Q[l,m,wd=25mm]Q[l,m,wd=40mm]},
    rowhead = 1,
    row{odd} = {bg=CinzaClaro},
    row{1} = {bg=AzulEscuro, fg=white},
    row{23} = {bg=AzulClaro, fg=white},
    cells  = {font=\fontsize{10pt}{12pt}\selectfont},
    }
        \textbf{Período} & \textbf{Componente} & \textbf{Código} & \textbf{Conteúdos}\\
        1\textordmasculine & Introdução à Prototipagem Eletrônica & FEELT!IPE & {\scriptsize circuitos elétricos; eletrônica digital; projeto de sistemas digitais; transdutores} \\
    2\textordmasculine & Física Básica: Mecânica & INFIS39031 & {\scriptsize física} \\
    2\textordmasculine & Experimental de Física Básica: Mecânica & INFIS39032 & {\scriptsize física} \\
    2\textordmasculine & Metrologia & FEELT31204 & {\scriptsize transdutores} \\
    2\textordmasculine & Sistemas Digitais & FEELT31409 & {\scriptsize eletrônica digital; projeto de sistemas digitais} \\
    2\textordmasculine & Experimental de Sistemas Digitais & FEELT31410 & {\scriptsize eletrônica digital; projeto de sistemas digitais} \\
    3\textordmasculine & Elementos de Sistemas Computacionais & FEELT!ESC & {\scriptsize eletrônica digital; microeletrônica e nanoeletrônica; projeto de circuitos integrados; projeto de sistemas digitais} \\
    3\textordmasculine & Física Básica: Eletricidade e Magnetismo & INFIS39033 & {\scriptsize eletricidade; física; teoria eletromagnética} \\
    3\textordmasculine & Experimental de Física Básica: Eletricidade e Magnetismo & INFIS39034 & {\scriptsize eletricidade; física; teoria eletromagnética} \\
    4\textordmasculine & Circuitos Elétricos I & FEELT31301 & {\scriptsize circuitos elétricos; eletricidade} \\
    4\textordmasculine & Experimental de Circuitos Elétricos I & FEELT31302 & {\scriptsize circuitos elétricos; eletricidade} \\
    5\textordmasculine & Eletrônica Analógica I & FEELT31401 & {\scriptsize circuitos elétricos; eletrônica analógica; microeletrônica e nanoeletrônica; teoria dos semicondutores} \\
    5\textordmasculine & Experimental de Eletrônica Analógica I & FEELT31402 & {\scriptsize circuitos elétricos; eletrônica analógica; microeletrônica e nanoeletrônica; teoria dos semicondutores} \\
    6\textordmasculine & Processamento Digital de Sinais & FEELT31628 & {\scriptsize processamento digital de sinais} \\
    6\textordmasculine & Redes de Comunicações I & FEELT31724 & {\scriptsize comunicação de dados} \\
    6\textordmasculine & Sistemas Embarcados I & FEELT31523 & {\scriptsize automação de projeto; eletrônica digital; microeletrônica e nanoeletrônica; projeto de sistemas digitais; sistemas de controle; transdutores} \\
    7\textordmasculine & Redes de Comunicações II & FEELT31831 & {\scriptsize comunicação de dados} \\
    7\textordmasculine & Sistemas Computacionais em Tempo Real & FEELT!SCTR & {\scriptsize automação de projeto; projeto de sistemas digitais; sistemas de controle} \\
    7\textordmasculine & Sistemas e Controle & FEELT!SCON & {\scriptsize automação de projeto; sistemas de controle} \\
    8\textordmasculine & Sistemas Distribuídos & FEELT!SD & {\scriptsize comunicação de dados} \\
    Optativa & Projetos Digitais com FPGA & FEELT!FPGA & {\scriptsize projeto de circuitos integrados} \\*

    \end{longtblr}
    
